% Generálva: html/build_tex.js (ai10.js). Kézzel ne szerkeszd, a forrás a html/data mappa.
\setcounter{chapter}{9}
\chapter{\texorpdfstring{Cellás mobilhálózat (GSM)}{Cellás mobilhálózat (GSM)}}\label{f10}
Kapacitásbecslés klaszterekkel és frekvencia-újrafelhasználással, valamint a roamingos „melyik HLR, melyik SMSC” kérdések. Ezekhez néhány szabályt kell biztosan tudni, a válaszoknál mindig ott van, melyikből jön.

\section{Alapok, egyszerűen}
\subsection{\texorpdfstring{A1. Kapacitás: vivők, időrések, klaszter}{A1. Kapacitás: vivők, időrések, klaszter}}\label{a10-kapacitas}
GSM 900: irányonként 25 MHz, 200 kHz-es vivők (124 db), vivőnként 8 időrés (TDMA). Egy \textbf{klaszter} $N$ cellából áll ($N$ = FRF, frekvencia-újrafelhasználási tényező): a vivőket szétosztják köztük, így egy klaszter minden vivőt pontosan egyszer használ. $C$ vivővel egy cellában $\approx C/N$ vivő, egyszerre $8C/N$ beszélgetés (a jelzési időréseket elhanyagolva), klaszterenként $8C$.

Nagy $N$: az azonos frekvenciájú cellák távolabb vannak, kisebb az interferencia, de kevesebb vivő jut egy cellára.

\subsection{\texorpdfstring{A2. Elemek és szabályok}{A2. Elemek és szabályok}}\label{a10-roaming}
MS: készülék. BTS: bázisállomás (rádió). BSC: bázisállomás-vezérlő. MSC: kapcsolóközpont. GMSC: átjáró MSC, a kívülről jövő hívást fogadja és a HLR-t kérdezi. HLR: honi adatbázis (hol van az előfizető). VLR: látogatói adatbázis. AuC: hitelesítés. EIR: készülék-nyilvántartás. SMSC: SMS-központ. Az index ($\alpha$, $\beta$) a szolgáltatót jelöli.

\begin{enumerate}
\item \textbf{Hívásnál a hívott honi HLR-jét} kérdezik le, a hívott honi hálózatának GMSC-je értesül először.
\item \textbf{Az SMS-t a küldő honi SMSC-je} fogadja és kézbesíti.
\item \textbf{Beszélgetés közbeni mozgásról} (hívásátadás) csak a kiszolgáló hálózat BTS-e, BSC-je és MSC-je tud. Hívás nélküli mozgásnál helyzetfrissítés van: a VLR, és VLR-váltáskor a honi HLR is értesül.
\item Egy eseményről (hívásindítás, bontás) először a kezdeményező saját BTS-e tud.
\end{enumerate}

Hitelesítés: AuC (Ki, A3 $\to$ SRES, A8 $\to$ Kc), rejtjelezés A5-tel. Készülékellenőrzés: a kiszolgáló hálózat EIR-je (fehér, fekete, szürke = megfigyelt lista). FNR: számhordozhatóság.

% ------------------------------------------------------------------
\section{\texorpdfstring{10.1. feladat: N = 7 klaszter, 70 csatorna (Moodle 2020)}{10.1. feladat: N = 7 klaszter, 70 csatorna (Moodle 2020)}}\label{ai10_01}
\setcounter{lepes}{0}
\begin{feladat}{10.1. feladat}
a) Rajzoljon N = 7-es klasztert két szomszédjával. b) 70 FDM-csatornával hány aktív előfizető lehet a klaszterben? c) Mi az FRF szerepe, mi a nagy N előnye és hátránya?

{\footnotesize\sffamily Forrás: ALL-IN 6. oldal $\cdot$ ai2020\_05\_26\_Q6. Címkék: vizsga}
\end{feladat}
\begin{keresunk}
\begin{itemize}
\item a) rajz
\item b) egyszerre aktív előfizetők a klaszterben
\item c) indoklás
\end{itemize}
\end{keresunk}
\begin{jeloles}
\begin{jeltabla}
$\ca{N}$ & a klaszter cellaszáma (FRF) & $7$\\
\end{jeltabla}
\par\smallskip\textbf{Adott még:}
\begin{itemize}
\item 70 FDM-csatorna (vivő), vivőnként 8 időrés
\end{itemize}
\end{jeloles}

\lepes{a) Rajz}
Egy középső hatszög és körülötte hat: ez egy klaszter, 1-től 7-ig számozva. Két szomszéd klaszter ugyanígy, eltolva, úgy, hogy az azonos számú cellák egymástól a lehető legmesszebb kerüljenek.

\lepes{b) Aktív előfizetők}
A klaszter minden vivőt egyszer használ fel (A1), minden vivőn 8 időrés.

\begin{center}\adjustbox{max width=\linewidth}{$\displaystyle 70\cdot8=560$}\end{center}
\jegyzet{számérték}{klaszterenként}

\lepes{c) Az FRF szerepe}
Megadja, hány cella után lehet ugyanazt a vivőt újra használni. Nagy N: az azonos vivőt használó cellák távolabb, kisebb az azonos csatornás interferencia. Hátrány: cellánként kevesebb vivő jut (pl. 120 vivő, N = 12: csak 10 vivő, 80 beszélgetés). Tipikus értékek: 3, 7, 12.

\begin{rajz}
\textbf{N = 7-es klaszter két szomszédjával}\par
\begin{center}\begin{tikzpicture}[x=0.02730cm,y=-0.02730cm,line cap=round,line join=round]
\useasboundingbox (-86,-79) rectangle (216,214);
\path[fill=black!10,draw=none] (26,-15) -- (26,15) -- (0,30) -- (-26,15) -- (-26,-15) -- (0,-30) -- cycle;
\node[anchor=base,inner sep=0pt,font=\footnotesize] at (0,5) {1};
\path[fill=black!10,draw=none] (77.9,-15) -- (77.9,15) -- (52,30) -- (26,15) -- (26,-15) -- (52,-30) -- cycle;
\node[anchor=base,inner sep=0pt,font=\footnotesize] at (52,5) {2};
\path[fill=black!10,draw=none] (52,-60) -- (52,-30) -- (26,-15) -- (0,-30) -- (0,-60) -- (26,-75) -- cycle;
\node[anchor=base,inner sep=0pt,font=\footnotesize] at (26,-40) {3};
\path[fill=black!10,draw=none] (0,-60) -- (0,-30) -- (-26,-15) -- (-52,-30) -- (-52,-60) -- (-26,-75) -- cycle;
\node[anchor=base,inner sep=0pt,font=\footnotesize] at (-26,-40) {4};
\path[fill=black!10,draw=none] (-26,-15) -- (-26,15) -- (-52,30) -- (-77.9,15) -- (-77.9,-15) -- (-52,-30) -- cycle;
\node[anchor=base,inner sep=0pt,font=\footnotesize] at (-52,5) {5};
\path[fill=black!10,draw=none] (0,30) -- (0,60) -- (-26,75) -- (-52,60) -- (-52,30) -- (-26,15) -- cycle;
\node[anchor=base,inner sep=0pt,font=\footnotesize] at (-26,50) {6};
\path[fill=black!10,draw=none] (52,30) -- (52,60) -- (26,75) -- (0,60) -- (0,30) -- (26,15) -- cycle;
\node[anchor=base,inner sep=0pt,font=\footnotesize] at (26,50) {7};
\path[fill=cA!15,draw=none] (155.9,30) -- (155.9,60) -- (129.9,75) -- (103.9,60) -- (103.9,30) -- (129.9,15) -- cycle;
\node[anchor=base,inner sep=0pt,font=\footnotesize] at (129.9,50) {1};
\path[fill=cA!15,draw=none] (207.8,30) -- (207.8,60) -- (181.9,75) -- (155.9,60) -- (155.9,30) -- (181.9,15) -- cycle;
\node[anchor=base,inner sep=0pt,font=\footnotesize] at (181.9,50) {2};
\path[fill=cA!15,draw=none] (181.9,-15) -- (181.9,15) -- (155.9,30) -- (129.9,15) -- (129.9,-15) -- (155.9,-30) -- cycle;
\node[anchor=base,inner sep=0pt,font=\footnotesize] at (155.9,5) {3};
\path[fill=cA!15,draw=none] (129.9,-15) -- (129.9,15) -- (103.9,30) -- (77.9,15) -- (77.9,-15) -- (103.9,-30) -- cycle;
\node[anchor=base,inner sep=0pt,font=\footnotesize] at (103.9,5) {4};
\path[fill=cA!15,draw=none] (103.9,30) -- (103.9,60) -- (77.9,75) -- (52,60) -- (52,30) -- (77.9,15) -- cycle;
\node[anchor=base,inner sep=0pt,font=\footnotesize] at (77.9,50) {5};
\path[fill=cA!15,draw=none] (129.9,75) -- (129.9,105) -- (103.9,120) -- (77.9,105) -- (77.9,75) -- (103.9,60) -- cycle;
\node[anchor=base,inner sep=0pt,font=\footnotesize] at (103.9,95) {6};
\path[fill=cA!15,draw=none] (181.9,75) -- (181.9,105) -- (155.9,120) -- (129.9,105) -- (129.9,75) -- (155.9,60) -- cycle;
\node[anchor=base,inner sep=0pt,font=\footnotesize] at (155.9,95) {7};
\path[fill=cB!15,draw=none] (52,120) -- (52,150) -- (26,165) -- (0,150) -- (0,120) -- (26,105) -- cycle;
\node[anchor=base,inner sep=0pt,font=\footnotesize] at (26,140) {1};
\path[fill=cB!15,draw=none] (103.9,120) -- (103.9,150) -- (77.9,165) -- (52,150) -- (52,120) -- (77.9,105) -- cycle;
\node[anchor=base,inner sep=0pt,font=\footnotesize] at (77.9,140) {2};
\path[fill=cB!15,draw=none] (77.9,75) -- (77.9,105) -- (52,120) -- (26,105) -- (26,75) -- (52,60) -- cycle;
\node[anchor=base,inner sep=0pt,font=\footnotesize] at (52,95) {3};
\path[fill=cB!15,draw=none] (26,75) -- (26,105) -- (0,120) -- (-26,105) -- (-26,75) -- (0,60) -- cycle;
\node[anchor=base,inner sep=0pt,font=\footnotesize] at (0,95) {4};
\path[fill=cB!15,draw=none] (0,120) -- (0,150) -- (-26,165) -- (-52,150) -- (-52,120) -- (-26,105) -- cycle;
\node[anchor=base,inner sep=0pt,font=\footnotesize] at (-26,140) {5};
\path[fill=cB!15,draw=none] (26,165) -- (26,195) -- (0,210) -- (-26,195) -- (-26,165) -- (0,150) -- cycle;
\node[anchor=base,inner sep=0pt,font=\footnotesize] at (0,185) {6};
\path[fill=cB!15,draw=none] (77.9,165) -- (77.9,195) -- (52,210) -- (26,195) -- (26,165) -- (52,150) -- cycle;
\node[anchor=base,inner sep=0pt,font=\footnotesize] at (52,185) {7};
\end{tikzpicture}\end{center}
\begin{itemize}
\item Egy klaszter: egy középső hatszög és a hat szomszédja, 1-től 7-ig számozva. A szomszéd klasztereket úgy told el (két cella egy irányba, egy cella 60\textdegree{}-kal elfordulva), hogy az azonos számú cellák egymástól a lehető legmesszebb legyenek.
\end{itemize}
\end{rajz}
\begin{eredmeny}

\begin{center}\adjustbox{max width=\linewidth}{$\displaystyle \text{b) }560\ \text{aktív előfizető}$}\end{center}
\end{eredmeny}


% ------------------------------------------------------------------
\section{\texorpdfstring{10.2. feladat: Győr: klaszter, egyidejű hívások, VAD, FNR}{10.2. feladat: Győr: klaszter, egyidejű hívások, VAD, FNR}}\label{ai10_02}
\setcounter{lepes}{0}
\begin{feladat}{10.2. feladat}
Győr 170 km\textsuperscript{2}, FRF = 12, kb. 85 GSM-cella. a) Mekkora egy klaszter? b) Hányan telefonálhatnak egyszerre? c) Egy szolgáltató 40 FDM-csatornát használ: hányan egy cellában nála? d) Mi a VAD előnye? e) Mire az FNR és az AuC?

{\footnotesize\sffamily Forrás: ALL-IN 25\ldots{}27., 145. oldal $\cdot$ ai2016\_06\_07\_3. Címkék: vizsga}
\end{feladat}
\begin{keresunk}
\begin{itemize}
\item a) a klaszter területe
\item b) egyidejű hívások
\item c) egy cellában egy szolgáltatónál
\item d), e) fogalmak
\end{itemize}
\end{keresunk}
\begin{jeloles}
\begin{jeltabla}
$T_\text{cella}$ & egy cella területe & \\
$T_\text{klaszter}$ & egy klaszter területe & \\
\end{jeltabla}
\par\smallskip\textbf{Adott még:}
\begin{itemize}
\item 124 vivő, vivőnként 8 időrés
\item c) egy szolgáltató 40 vivőt használ
\end{itemize}
\end{jeloles}

\lepes{a) Klaszter}
Először egy cella területe, majd 12 cella.

\begin{center}\adjustbox{max width=\linewidth}{$\displaystyle T_\text{cella}=\frac{170}{85}=2\ \text{km}^2\qquad T_\text{klaszter}=12\cdot2=24\ \text{km}^2$}\end{center}
\jegyzet{számérték}{}

\lepes{b) Egyidejű hívások}
Egy cellára $124\cdot8/12=82{,}7$ beszélgetés jut (A1).

\begin{center}\adjustbox{max width=\linewidth}{$\displaystyle 85\cdot\frac{124\cdot8}{12}=85\cdot82{,}7\approx7000$}\end{center}
\jegyzet{számérték}{ugyanez: $85/12\approx7$ klaszter $\cdot992$}

\lepes{c) Egy szolgáltató}
A 40 vivő oszlik meg a klaszter 12 cellája között.

\begin{center}\adjustbox{max width=\linewidth}{$\displaystyle \frac{40}{12}=3{,}33\ \text{vivő}\ \Rightarrow\ 3\cdot8=24\ \text{(néhol }4\cdot8=32)$}\end{center}
\jegyzet{számérték}{átlagosan kb. 26}

\lepes{d), e) Fogalmak}
VAD + DTX: beszédszünetben nincs adás, kevesebb fogyasztás és interferencia, a vevő komfortzajt generál. FNR: számhordozhatósági adatbázis. AuC: a SIM titkos kulcsának (Ki) tárolása, a hitelesítési adatok (RAND, SRES, Kc) előállítása.
\begin{eredmeny}

\begin{center}\adjustbox{max width=\linewidth}{$\displaystyle \text{a) }24\ \text{km}^2\qquad\text{b) }\approx7000\qquad\text{c) }\approx24\ldots32$}\end{center}
\end{eredmeny}
\begin{tipp}[Ellenőrzés]
\begin{itemize}
\item c) A gyűjteményben „passz”, a saját megoldás feltevése, hogy a 40 vivő a klaszterre vonatkozik \texttt{[ELLENŐRIZENDŐ]}.
\end{itemize}
\end{tipp}


% ------------------------------------------------------------------
\section{\texorpdfstring{10.3. feladat: Két roamingoló előfizető hívása és SMS-e}{10.3. feladat: Két roamingoló előfizető hívása és SMS-e}}\label{ai10_03}
\setcounter{lepes}{0}
\begin{feladat}{10.3. feladat}
MS-A az $\alpha$, MS-B a $\beta$ előfizetője, egymás hálózatában barangolnak. MS-A felhívja MS-B-t. a) Melyik HLR-t kérdezik? b) Melyik GMSC értesül előbb? c) Ki tud MS-A beszélgetés közbeni mozgásáról? d) Bontás után MS-B SMS-t küld MS-A-nak, melyik SMSC kézbesíti?

{\footnotesize\sffamily Forrás: ALL-IN 21\ldots{}22., 41. oldal $\cdot$ ai2016\_01\_05\_7, ai2016\_01\_05\_7b (és 5-ösöm 27. oldal). Címkék: vizsga, javítás}
\end{feladat}
\begin{keresunk}
\begin{itemize}
\item a) HLR
\item b) GMSC
\item c) a mozgásról tudó elemek
\item d) SMSC
\end{itemize}
\end{keresunk}
\begin{jeloles}
\begin{jeltabla}
$\beta$ & a másik szolgáltató ($\beta$) hálózata, alsó indexben: a $\beta$ hálózat eleme & \\
$\alpha$ & az egyik szolgáltató ($\alpha$) hálózata, alsó indexben: az $\alpha$ hálózat eleme & \\
\end{jeltabla}
\par\smallskip\textbf{Adott még:}
\begin{itemize}
\item MS-A: $\alpha$-előfizető, $\beta$-ban van
\item MS-B: $\beta$-előfizető, $\alpha$-ban van
\end{itemize}
\end{jeloles}

\lepes{A hívás útja}
MS-A $\beta$-ban hív: BTS$_\beta$ $\to$ BSC$_\beta$ $\to$ MSC$_\beta$. A hívott szám $\beta$-beli, ezért a GMSC$_\beta$ fogadja, megkérdezi a HLR$_\beta$-t, amely tudja, hogy MS-B a VLR$_\alpha$ területén van. A hívás innen az $\alpha$ hálózatba megy MS-B-hez.

\lepes{a), b)}
1. szabály: a hívott (MS-B, $\beta$-előfizető) honi HLR-je, és a honi hálózatának GMSC-je.

\begin{center}\adjustbox{max width=\linewidth}{$\displaystyle \text{HLR}_\beta,\ \text{GMSC}_\beta$}\end{center}
\jegyzet{definíció}{}

\lepes{c) Mozgás beszélgetés közben}
3. szabály: hívásátadás MS-A kiszolgáló hálózatában ($\beta$).

\begin{center}\adjustbox{max width=\linewidth}{$\displaystyle \text{BTS}_\beta,\ \text{BSC}_\beta,\ \text{MSC}_\beta$}\end{center}
\jegyzet{definíció}{}

\lepes{d) SMS}
2. szabály: a küldő (MS-B, $\beta$-előfizető) honi SMSC-je.

\begin{center}\adjustbox{max width=\linewidth}{$\displaystyle \text{SMSC}_\beta$}\end{center}
\jegyzet{definíció}{}
\begin{figyelem}[Javítás]
A c) részre a gyűjtemény a HLR$_\alpha$-t és a VLR$_\beta$-t is felsorolja. Hívás közben nincs helyzetfrissítés, ezek csak hívás nélküli mozgásnál értesülnek \texttt{[ELLENŐRIZENDŐ]}.
\end{figyelem}
\begin{eredmeny}

\begin{center}\adjustbox{max width=\linewidth}{$\displaystyle \text{a) HLR}_\beta\quad\text{b) GMSC}_\beta\quad\text{c) BTS}_\beta,\ \text{BSC}_\beta,\ \text{MSC}_\beta\quad\text{d) SMSC}_\beta$}\end{center}
\end{eredmeny}


% ------------------------------------------------------------------
\section{\texorpdfstring{10.4. feladat: Két $\alpha$-előfizető, az egyik $\beta$-ban (Moodle 2020)}{10.4. feladat: Két α-előfizető, az egyik β-ban (Moodle 2020)}}\label{ai10_04}
\setcounter{lepes}{0}
\begin{feladat}{10.4. feladat}
MS-A1 és MS-A2 is $\alpha$-előfizető, MS-A2 a $\beta$ hálózatban barangol. MS-A1 hívja MS-A2-t. a) Melyik eszközhöz érkezik először információ a hívásindításról? b) Melyik HLR-t kérdezik? c) Mely eszközök tudnak MS-A2 $\beta$-beli mozgásáról? d) Melyik eszközzel azonosítják MS-A1 SIM-jét?

{\footnotesize\sffamily Forrás: ALL-IN 8. oldal $\cdot$ ai2020\_05\_26\_Q8 (és 5-ösöm 28. oldal). Címkék: vizsga}
\end{feladat}
\begin{keresunk}
\begin{itemize}
\item a) az első eszköz
\item b) HLR
\item c) a mozgásról tudó elemek
\item d) a SIM azonosítása
\end{itemize}
\end{keresunk}
\begin{jeloles}
\begin{jeltabla}
$\alpha$ & az egyik szolgáltató ($\alpha$) hálózata, alsó indexben: az $\alpha$ hálózat eleme & \\
$\beta$ & a másik szolgáltató ($\beta$) hálózata, alsó indexben: a $\beta$ hálózat eleme & \\
\end{jeltabla}
\par\smallskip\textbf{Adott még:}
\begin{itemize}
\item MS-A1: $\alpha$-előfizető, $\alpha$-ban
\item MS-A2: $\alpha$-előfizető, $\beta$-ban
\end{itemize}
\end{jeloles}

\lepes{a) Hívásindítás}
4. szabály: a kezdeményező (MS-A1) saját BTS-e.

\begin{center}\adjustbox{max width=\linewidth}{$\displaystyle \text{BTS}_\alpha$}\end{center}
\jegyzet{definíció}{}

\lepes{b) HLR}
1. szabály: a hívott (MS-A2) $\alpha$-előfizető.

\begin{center}\adjustbox{max width=\linewidth}{$\displaystyle \text{HLR}_\alpha$}\end{center}
\jegyzet{definíció}{ez tudja, hogy MS-A2 a VLR$_\beta$-nál van}

\lepes{c) Mozgás $\beta$-ban}
Itt a mozgás általában (nem beszélgetés közben) értendő: a helyzetfrissítés elemei is számítanak. A rádiós elemek (BTS, BSC) és az MSC/VLR a kiszolgáló hálózatban, VLR-váltáskor a honi HLR is.

\begin{center}\adjustbox{max width=\linewidth}{$\displaystyle \text{BTS}_\beta,\ \text{BSC}_\beta,\ \text{MSC}_\beta,\ \text{VLR}_\beta,\ \text{HLR}_\alpha$}\end{center}
\jegyzet{definíció}{}

\lepes{d) SIM azonosítása}
A SIM titkos kulcsát (Ki) csak a honi AuC ismeri, ez állítja elő a hitelesítési adatokat.

\begin{center}\adjustbox{max width=\linewidth}{$\displaystyle \text{AuC}_\alpha$}\end{center}
\jegyzet{definíció}{}
\begin{eredmeny}

\begin{center}\adjustbox{max width=\linewidth}{$\displaystyle \text{a) BTS}_\alpha\quad\text{b) HLR}_\alpha\quad\text{c) BTS}_\beta,\ \text{BSC}_\beta,\ \text{MSC}_\beta,\ \text{VLR}_\beta,\ \text{HLR}_\alpha\quad\text{d) AuC}_\alpha$}\end{center}
\end{eredmeny}
\begin{tipp}[Ellenőrzés]
\begin{itemize}
\item A hallgató c)-ből kihagyta a BTS-t és a BSC-t, a javító ezért vont le pontot.
\end{itemize}
\end{tipp}


% ------------------------------------------------------------------
\section{\texorpdfstring{10.5. feladat: GSM-hitelesítés: A3, A8, A5}{10.5. feladat: GSM-hitelesítés: A3, A8, A5}}\label{ai10_05}
\setcounter{lepes}{0}
\begin{feladat}{10.5. feladat}
Az ábrán AuC/HLR, MSC/VLR és MS. Milyen adatokat adnak egymásnak (A, B, C)? Mik a szürke dobozok? Mire való a Kc?

{\footnotesize\sffamily Forrás: ALL-IN 87\ldots{}88. oldal $\cdot$ ai2016\_01\_12c\_5. Címkék: vizsga}
\end{feladat}
\begin{keresunk}
\begin{itemize}
\item A, B, C, a dobozok, Kc
\end{itemize}
\end{keresunk}
\begin{jeloles}
\begin{jeltabla}
$A$ & az ábra A jelű üzenete & \\
$B$ & az ábra B jelű üzenete & \\
$C$ & az ábra C jelű üzenete & \\
\end{jeltabla}
\par\smallskip\textbf{Adott még:}
\begin{itemize}
\item Ki: a SIM titkos kulcsa, csak a SIM és az AuC ismeri
\end{itemize}
\end{jeloles}

\lepes{Hitelesítés}
Az AuC véletlen számot (RAND, 128 bit) állít elő, ebből és a Ki-ből az A3 algoritmussal a várt választ (SRES, 32 bit). A RAND-ot és az SRES-t az MSC/VLR kapja meg. A RAND a készülékhez megy (A), a SIM ugyanazzal az A3-mal kiszámolja a saját SRES-ét és visszaküldi (C). Az MSC összeveti az AuC-tól kapottal (B): ha egyezik, a SIM valódi.

\lepes{Rejtjelezés}
Az A8 algoritmus mindkét oldalon a RAND-ból és a Ki-ből előállítja a rejtjelező kulcsot (Kc, 64 bit). Az A5 ebből rejtjelező bitsorozatot képez, ezt a rádiós burstök 114 adatbitjével XOR-olják. A Kc soha nem megy át a rádiós szakaszon.
\begin{eredmeny}

\begin{center}\adjustbox{max width=\linewidth}{$\displaystyle A=\text{RAND},\ B=\text{SRES (AuC-tól)},\ C=\text{SRES (MS-től)}\qquad\text{dobozok: A3, A8 és A5}\qquad\text{Kc: a rádiós szakasz titkosítása}$}\end{center}
\end{eredmeny}


% ------------------------------------------------------------------
\section{\texorpdfstring{10.6. feladat: VAD, EIR, FNR, AuC röviden}{10.6. feladat: VAD, EIR, FNR, AuC röviden}}\label{ai10_06}
\setcounter{lepes}{0}
\begin{feladat}{10.6. feladat}
Hogyan működik a VAD, milyen erőforrás takarítható meg vele (GSM, VoIP)? Mi az EIR, az FNR és az AuC szerepe?

{\footnotesize\sffamily Forrás: ALL-IN 27., 54\ldots{}55., 82. oldal $\cdot$ ai2016\_06\_07\_3, ai2014\_12\_15\_3, ai2014\_01\_14\_3. Címkék: beugró, javítás}
\end{feladat}
\begin{keresunk}
\begin{itemize}
\item indokolt válaszok
\end{itemize}
\end{keresunk}
\begin{jeloles}
\par\smallskip\textbf{Adott még:}
\begin{itemize}
\item fogalmi kérdések
\end{itemize}
\end{jeloles}

\lepes{VAD}
A hangaktivitás-érzékelő felismeri a beszédszünetet, ilyenkor a DTX leállítja az adást. GSM-ben az akkumulátor és a rádiós interferencia, VoIP-ban a sávszélesség (kevesebb csomag) takarítható meg. A vevő komfortzajt generál, hogy a vonal ne tűnjön megszakadtnak.

\lepes{EIR, FNR, AuC}
EIR: IMEI-listák. Fehér: engedélyezett, fekete: tiltott (pl. lopott), szürke: megfigyelt. FNR: számhordozhatósági adatbázis, megmondja, melyik szolgáltatóhoz került egy hívószám. AuC: a Ki tárolása és a hitelesítési adatok (RAND, SRES, Kc) előállítása.
\begin{figyelem}[Javítás]
A gyűjtemény szerint a szürke listás készülékek „minden lépése blokkolva van”: helyesen használhatók, csak megfigyeltek.
\end{figyelem}
\begin{eredmeny}

\begin{center}\adjustbox{max width=\linewidth}{$\displaystyle \text{VAD: szünetben nincs adás}\quad\text{EIR: IMEI-listák}\quad\text{FNR: számhordozhatóság}\quad\text{AuC: hitelesítés}$}\end{center}
\end{eredmeny}
